\documentclass[11pt]{article}
\usepackage[letterpaper,margin=0.72in,top=0.82in,bottom=0.72in,headheight=22pt]{geometry}
\usepackage[T1]{fontenc}
\usepackage{lmodern}
\usepackage{microtype}
\usepackage[table]{xcolor}
\usepackage{amsmath,amssymb,mathtools}
\usepackage{array,longtable,booktabs,tabularx}
\usepackage[inline]{enumitem}
\usepackage[most]{tcolorbox}
\usepackage{fancyhdr}
\usepackage{hyperref}

\definecolor{Primary}{HTML}{17324D}
\definecolor{Accent}{HTML}{2C7A7B}
\definecolor{CueTint}{HTML}{EAF1F5}
\definecolor{NoteTint}{HTML}{F8FAFB}
\definecolor{Rule}{HTML}{B8C6CF}
\hypersetup{hidelinks,pdftitle={Chapter 16: Maps, Mapping, Sequencing, and Superstrings},pdfsubject={Cornell notes}}
\setlist[itemize]{leftmargin=1.25em,itemsep=1pt,topsep=1.5pt,parsep=0pt}
\setlength{\parindent}{0pt}
\setlength{\parskip}{3pt}
\renewcommand{\arraystretch}{1.16}
\emergencystretch=2em

\pagestyle{fancy}
\fancyhf{}
\fancyhead[L]{\sffamily\small\color{Primary} Chapter 16}
\fancyhead[R]{\sffamily\small\color{Primary} Mapping and superstrings}
\fancyfoot[C]{\sffamily\small\color{Primary}\thepage}
\renewcommand{\headrulewidth}{0.5pt}
\renewcommand{\headrule}{\hbox to\headwidth{\color{Accent}\leaders\hrule height \headrulewidth\hfill}}

\newcolumntype{C}{>{\raggedright\arraybackslash}p{0.275\textwidth}}
\newcolumntype{N}{>{\raggedright\arraybackslash}p{0.655\textwidth}}
\newcommand{\cnrow}[2]{%
  \cellcolor{CueTint}\textbf{\sffamily\color{Primary}#1} &
  \cellcolor{NoteTint}#2 \\[4pt]\arrayrulecolor{Rule}\hline}

\begin{document}

\begin{tcolorbox}[
  enhanced,breakable,colback=Primary!4,colframe=Primary,boxrule=0.9pt,
  arc=2mm,left=3mm,right=3mm,top=2.5mm,bottom=2.5mm]
{\sffamily\color{Accent}\bfseries CHAPTER 16}\\[-1pt]
{\sffamily\LARGE\bfseries\color{Primary} Maps, Mapping, Sequencing, and Superstrings}

\vspace{2pt}
\textbf{Purpose.} Translate genome mapping and sequencing workflows into ordering, overlap, layout, assembly, and shortest-superstring problems.

\textbf{Learning objectives}
\begin{itemize}
  \item Distinguish genetic and physical maps.
  \item Model STS, radiation-hybrid, and fingerprint mapping data.
  \item Compare directed, hierarchical, and shotgun sequencing strategies.
  \item Relate overlap assembly to shortest common superstrings and hybridization data.
\end{itemize}

\textbf{Prerequisites.} Graphs, intervals, permutations, probability, exact/approximate matching, and basic genomics.
\end{tcolorbox}

\begin{longtable}{@{}C!{\color{Accent}\vrule width 0.8pt}N@{}}
\rowcolor{Primary}
\color{white}\sffamily\bfseries Cue / question &
\color{white}\sffamily\bfseries Notes \\ \hline
\endfirsthead
\rowcolor{Primary}
\color{white}\sffamily\bfseries Cue / question &
\color{white}\sffamily\bfseries Notes (continued) \\ \hline
\endhead
\cnrow{\S16.1: What computational tasks recur?}{Mapping and sequencing repeatedly ask which fragments overlap, in what order they occur, how far apart landmarks lie, and which global layout best explains noisy local observations. Inputs are indirect measurements; algorithms must distinguish exact consistency from best-fit reconstruction.}
\cnrow{\S16.2: Why build a map?}{A map supplies an ordered scaffold of markers or clones before every base is known. It supports navigation, targeted sequencing, assembly validation, and localization of features. Resolution and error depend on the measurement technology.}
\cnrow{\S16.3: Genetic versus physical maps}{A genetic map estimates relative positions from recombination frequencies and is measured in recombination units. A physical map orders landmarks along DNA and seeks base-pair or clone-scale distance. Recombination rate varies, so genetic distance is not a uniform conversion to physical distance.}
\cnrow{\S16.4: What is physical mapping?}{Represent experimentally observed clones, probes, markers, or fragment overlaps and infer a consistent order and spacing. Repeats, missing observations, chimeric clones, and measurement noise turn an ideal interval-ordering problem into an optimization problem.}
\cnrow{\S16.5: How does STS-content mapping work?}{An STS is a uniquely detectable landmark. Record a binary matrix indicating which clone contains which STS. If clones are intervals on one chromosome, seek an STS order that makes the positive entries of every clone consecutive, then derive an ordered clone library and overlaps.}
\cnrow{What is the ideal consistency test?}{The consecutive-ones property captures whether every clone's marker set can be contiguous in one order. Specialized order representations can maintain all feasible orders. Real errors require scoring violations or identifying suspect rows rather than declaring the entire data set impossible.}
\cnrow{\S16.6: What does radiation-hybrid mapping measure?}{Radiation breaks chromosomes into fragments; markers retained together across hybrid cells are likely close. Pairwise retention patterns estimate separation and ordering. A probabilistic likelihood model is more appropriate than treating every observation as exact.}
\cnrow{\S16.7: How does fingerprint mapping work?}{Digest or otherwise fingerprint clones and compare fragment patterns. Sufficient similarity suggests overlap. Build an overlap graph or weighted compatibility structure, suppress repetitive evidence, and seek a layout consistent with clone intervals.}
\cnrow{\S16.8: What is the tightest layout?}{Given overlap or containment constraints, minimize the span while preserving all required relations. Remove contained or duplicate fragments when the model permits, convert overlaps into relative offsets, and test that cycles do not imply contradictory placements.}
\cnrow{\S16.9: What limits physical maps?}{False overlaps, repeats, clone gaps, nonuniform coverage, and local ambiguity can produce several plausible layouts. Report connected components, confidence, and unresolved order rather than forcing a single total order.}
\cnrow{\S16.10: What is map alignment?}{Compare two ordered marker maps when markers may be missing, duplicated, reversed, or measured at different scales. Dynamic programming or chaining selects order-consistent correspondences and penalizes gaps or conflicting intervals.}
\cnrow{\S16.11: What is large-scale sequencing?}{Sequencing generates reads or subclones and reconstructs longer sequence through experimental design plus assembly. The computational problem changes with read length, coverage, error rate, paired constraints, and whether a map or clone hierarchy is available.}
\cnrow{\S16.12: What is directed sequencing?}{Choose each next experiment from the end of already known sequence, extending into adjacent unknown DNA. It reduces assembly ambiguity but requires sequential laboratory decisions and can progress slowly.}
\cnrow{\S16.13: What is top-down/bottom-up sequencing?}{First order large clones into a physical map, then sequence and assemble within selected clones, and finally merge the clone-level results. The hierarchy localizes difficult assembly but introduces mapping and clone-management overhead.}
\cnrow{\S16.14: What is shotgun sequencing?}{Sample many reads without directing each start. Coverage and redundancy create overlaps from which assembly is inferred. Parallel data generation is fast, but repeats and uneven coverage create ambiguous graph structures.}
\cnrow{\S16.15: What is sequence assembly?}{Filter or correct reads, detect overlaps or shared words, construct an overlap graph or related layout structure, resolve paths with coverage and pairing evidence, and output contigs plus unresolved gaps. A valid assembly should also be checked against the reads and independent mapping evidence.}
\cnrow{What is the overlap-layout profile?}{\textbf{Input:} reads and approximate overlaps. \textbf{State:} vertices for reads or sequence contexts and weighted compatibility edges. \textbf{Goal:} a consistent path or layout covering evidence. \textbf{Risks:} repeats, contained reads, orientation, and sequencing errors. \textbf{Cost:} overlap discovery often dominates naive all-pairs comparison.}
\cnrow{\S16.16: How should strategies be compared?}{Directed approaches spend experimental coordination to simplify assembly; shotgun approaches spend coverage and computation; hierarchical approaches use maps to partition the problem. Hybrid strategies trade among these resources.}
\cnrow{\S16.17: What is shortest common superstring?}{Given strings, find the shortest string containing each as a substring. Remove contained strings, define directed overlap weights, and seek an order maximizing total reused overlap. The general optimization is computationally hard, motivating approximation and application-specific constraints.}
\cnrow{How does overlap relate to length?}{For an ordering $s_{\pi_1},\ldots,s_{\pi_k}$ without contained strings, the constructed superstring length equals total input length minus the sum of consecutive overlaps. Maximizing overlap is therefore the ordering objective, but greedy locally largest overlap need not be globally optimal.}
\cnrow{\S16.18: What is sequencing by hybridization?}{Infer a target string from the set or multiset of observed length-$\ell$ substrings. In an ideal spectrum, connect compatible $(\ell-1)$ contexts and reconstruct a path using every observation. Missing, duplicate, and spurious probes turn exact reconstruction into a noisy path or superstring problem.}
\cnrow{\S16.19: What should practice emphasize?}{Translate each experimental data type into its combinatorial object, identify orientation and noise assumptions, and separate feasibility, optimization, and validation. Construct counterexamples to greedy overlap assembly.}
\end{longtable}
\enlargethispage{1\baselineskip}

\begin{tcolorbox}[enhanced,breakable,title={Chapter synthesis},
  colback=Accent!5,colframe=Accent,fonttitle=\small\sffamily\bfseries,fontupper=\footnotesize,
  boxsep=0.6mm,left=1.5mm,right=1.5mm,top=1mm,bottom=1mm,before skip=3pt,after skip=3pt]
Mapping and sequencing infer global order from partial local evidence. Marker-content data suggest interval orders, hybrid retention suggests probabilistic distances, fingerprints and reads suggest overlaps, and assembly searches for a consistent layout. Shortest-superstring and spectrum formulations expose the combinatorial core but idealize away much of the noise handled by real workflows.
\end{tcolorbox}

\begin{tcolorbox}[enhanced,breakable,title={Key takeaways},
  colback=white,colframe=Primary,fonttitle=\small\sffamily\bfseries,fontupper=\footnotesize,
  boxsep=0.6mm,left=1.5mm,right=1.5mm,top=1mm,bottom=1mm,before skip=3pt,after skip=3pt]
\begin{itemize}
  \item Genetic distance and physical distance measure different phenomena.
  \item STS-content mapping seeks a marker order with contiguous clone intervals.
  \item Radiation-hybrid data are probabilistic.
  \item Fingerprint and read overlap graphs require repeat-aware validation.
  \item Directed, hierarchical, and shotgun strategies shift cost between laboratory work and computation.
  \item Shortest superstring maximizes reused overlap but is computationally hard.
  \item Hybridization spectra can be reconstructed through compatible context paths.
\end{itemize}
\end{tcolorbox}

\begin{tcolorbox}[enhanced,breakable,title={Algorithm selection / when to use this},
  colback=Primary!3,colframe=Primary,fonttitle=\small\sffamily\bfseries,fontupper=\footnotesize,
  boxsep=0.6mm,left=1.5mm,right=1.5mm,top=1mm,bottom=1mm,before skip=3pt,after skip=3pt]
\begin{itemize}
  \item Use consecutive-ones reasoning for ideal marker-content maps.
  \item Use likelihood-based ordering for noisy retention data.
  \item Use overlap or context graphs for read assembly and spectra.
  \item Use directed sequencing when controlled extension is worth sequential effort.
  \item Use shotgun sequencing when parallel throughput and coverage are available.
\end{itemize}
\end{tcolorbox}

\begin{tcolorbox}[enhanced,breakable,title={Self-test questions},
  colback=white,colframe=Accent,fonttitle=\small\sffamily\bfseries,fontupper=\footnotesize,
  boxsep=0.6mm,left=1.5mm,right=1.5mm,top=1mm,bottom=1mm,before skip=3pt,after skip=3pt]
\begin{itemize}
  \item Why can recombination distance fail to reflect physical distance uniformly?
  \item What matrix property characterizes an ideal STS-content map?
  \item How do repeats create false overlap edges?
  \item Compare directed, hierarchical, and shotgun sequencing.
  \item Why does maximizing consecutive overlap minimize an ordered superstring?
\end{itemize}
\end{tcolorbox}

{\footnotesize
\textbf{Terms and notation to remember.} genetic map, physical map, STS, radiation hybrid, fingerprint, overlap graph, contig, superstring, spectrum.

\textbf{Connections.} Indexes accelerate overlaps; rearrangement models compare large-scale order.
}

\end{document}
